\documentclass[10pt]{article}
\usepackage[T1]{fontenc}
\usepackage[utf8]{inputenc}
\usepackage[french]{babel}
\usepackage{amsmath,amssymb,mathrsfs,mathpazo}
\usepackage{geometry}
\geometry{a4paper,margin=2cm}
\usepackage{graphicx,array,tabularx,tabularray,colortbl,multirow,multicol,booktabs,enumitem,pifont,ulem,fancybox,diagbox,cancel}
\usepackage{pstricks,pstricks-add,pst-plot,pst-tree,pst-node,pst-text,pst-func,pst-math,pst-bspline,pst-3dplot,pst-eucl,pst-solides3d}
\let\pscalculate\undefined
\usepackage{pst-bezier}
\usepackage{tikz,pgfplots,tkz-tab}
\usetikzlibrary{arrows,arrows.meta,patterns,calc,positioning,trees,decorations.markings,shapes.geometric,shapes.misc,shapes.arrows,intersections,decorations.pathmorphing,decorations.pathreplacing,backgrounds,shadows,babel}
\pgfplotsset{compat=1.18}
\usepgfplotslibrary{fillbetween}
\usepackage[np]{numprint}
\usepackage{listings,makecell,pas-tableur,twoopt}
\let\barre\undefined
\usepackage[pstricks]{tablvar}
\definecolor{darkgreen}{RGB}{0,100,0}
\providecommand{\red}{\color{red}}
\providecommand{\blue}{\color{blue}}
\providecommand{\green}{\color{green}}
\providecommand{\black}{\color{black}}
\providecommand{\white}{\color{white}}
\providecommand{\gray}{\color{gray}}
\providecommand{\cyan}{\color{cyan}}
\providecommand{\magenta}{\color{magenta}}
\providecommand{\yellow}{\color{yellow}}
\providecommand{\violet}{\color{violet}}
\providecommand{\orange}{\color{orange}}
\providecommand{\euro}{\texteuro}
\providecommand{\EUR}{\texteuro}
\providecommand{\ieme}{\textsuperscript{e}}
\providecommand{\degres}{\ensuremath{^\circ}}
\providecommand{\degree}{\ensuremath{^\circ}}
\providecommand{\vv}[1]{\overrightarrow{#1}}
\providecommand{\Coord}[1]{#1}
\providecommand{\decosix}{}
\providecommand{\starredbullet}{\ensuremath{\bullet}}
\pagestyle{empty}
\setlength{\parindent}{0pt}
\usepackage{adjustbox}\DeclareUnicodeCharacter{2236}{:}\DeclareUnicodeCharacter{2190}{\ensuremath{\leftarrow}}
\providecommand{\R}{\mathbb{R}}
\providecommand{\N}{\mathbb{N}}
\providecommand{\Z}{\mathbb{Z}}
\providecommand{\Q}{\mathbb{Q}}
\providecommand{\D}{\mathbb{D}}
\providecommand{\C}{\mathbb{C}}
\providecommand{\e}{{\mathrm{e}}}
\providecommand{\vect}[1]{\overrightarrow{#1}}
\providecommand{\vectt}[1]{\overrightarrow{\mathrm{#1}}}
\providecommand{\vvt}[1]{\overrightarrow{\mathrm{#1}}}
\providecommand{\np}[2][]{{#2~#1}}
\providecommand{\vv}[1]{{\overrightarrow{#1}}}
\renewcommand{\d}{{\mathrm{d}}}
\providecommand{\degres}{{^\circ}}
\providecommand{\degre}{{^\circ}}
\providecommand{\blue}{\color{blue}}
\providecommand{\red}{\color{red}}
\providecommand{\enskip}{\quad}
\providecommand{\llbracket}{\lbrack\!\lbrack}
\providecommand{\rrbracket}{\rbrack\!\rbrack}
\providecommand{\euro}{{\text{€}}}
\providecommand{\no}{{\text{n°}}}
\providecommand{\interval}[2]{{[#1;#2]}}
\providecommand{\intervFO}[2]{{[#1;#2[}}
\providecommand{\dotfill}{\ldots}

\providecommand{\ds}{\displaystyle}
\providecommand{\textding}[1]{#1}
\providecommand{\Oij}{\(\mathrm{O};\vec{\imath},\vec{\jmath}\)}
\providecommand{\Oijk}{\(\mathrm{O};\vec{\imath},\vec{\jmath},\vec{k}\)}
\providecommand{\Ouv}{\(\mathrm{O};\vec{u},\vec{v}\)}
\providecommand{\intervFF}[2]{\ensuremath{[#1;#2]}}
\renewcommand{\euro}{\ifmmode\text{\texteuro}\else\texteuro\fi}
\DeclareUnicodeCharacter{00B0}{\ifmmode^{\circ}\else\textdegree\fi}
\DeclareUnicodeCharacter{20AC}{\ifmmode\text{\texteuro}\else\texteuro\fi}
% Énoncés : Éducation nationale. Transcription et corrigés : APMEP, Denis Vergès et contributeurs. Découpage et classement : CoopMaths, Sébastien Lozano, Éric Elter et contributeurs. Adaptation de la présentation : Maths Explorer.
% CC BY-SA : https://coopmaths.fr/www/about/
% https://coopmaths.fr/alea/static/bac/2026/tex/bac_2026_09_sujet2_metropole_4_cor.tex
\begin{document}
\noindent\textit{APMEP / CoopMaths. CC BY-SA. Présentation adaptée par Maths Explorer.}\par\medskip
\subsection*{Partie A}

	\begin{enumerate}
		\item On a $\vectt{RS} \begin{pmatrix}
			0 \\ 3 \\ -4
		\end{pmatrix}$ et $\vectt{RT} \begin{pmatrix}
			\np{2,5} \\ 7 \\ -1
		\end{pmatrix}$

		Puisque $\dfrac{0}{\np{2,5}} \neq \dfrac{3}{7}$, les coordonnées des vecteurs $\vectt{RS}$ et $\vectt{RT}$ ne sont pas proportionnelles donc les vecteurs $\vectt{RS}$ et $\vectt{RT}$ ne sont pas colinéaires et les points R, S et T ne sont pas alignés.

		\item \begin{enumerate}
			\item Les vecteurs $\vectt{RS}$ et $\vectt{RT}$ ne sont pas colinéaires donc ils forment une base du plan (RST).

			$\vect{n}\cdot \vectt{RS} = 0+12-12=0$ \quad et \quad $\vect{n}\cdot \vectt{RT} = -25+28-3=0$.

			Le vecteur $\vect{n}$ est orthogonal à deux vecteurs formant une base du plan (RST) donc le vecteur $\vect{n}$ est normal au plan (RST).

			\item Le vecteur $\vect{n} \begin{pmatrix}-10 \\ 4 \\ 3\end{pmatrix} $ est normal au plan (RST) donc une équation du plan (RST) est de la forme \quad  $-10x+4y+3z+d=0$.

			Le point R$(1;0;6)$ appartient au plan (RST) donc on doit avoir \quad  $-10 \times 1 +4 \times 0 +3 \times 6 + d =0$\quad  soit \quad $d=-8$.

			Ainsi, l'équation \quad $-10x+4y+3z-8=0$\quad est une équation cartésienne du plan (RST).
		\end{enumerate}

		\item \begin{enumerate}
			\item On a \quad $-10x_{\mathrm{A}}+4y_{\mathrm{A}}+3z_{\mathrm{A}}-8 = -10 \times (\np{2,5}) +4 \times 8 +3 \times \np{4,5}  -8 =-25+32+\np{13,5}-8 \neq 0$.

			On en déduit que A n'appartient pas à (RST).

			\item Le vecteur $\vect{n} \begin{pmatrix}-10 \\ 4 \\ 3\end{pmatrix} $ est normal au plan (RST) donc c'est un vecteur directeur de la droite $d_1$.

			De plus, A$(\np{-2,5} ; 8 ; \np{4,5}) \in d_1$ donc une représentation paramétrique de $d_1$ est

			$\begin{cases}
				x=\np{-2,5} -10t \\
				y=8 + 4t \\
				z=\np{4,5} +3t
			\end{cases}, t \in \R$

			\item A $\in d_1$ et $d_1$ est orthogonale au plan (RST) donc le projeté orthogonal de A sur (RST) est le point d'intersection de $d_1$ et de (RST).

			Puisque\quad  $-10 \times (\np{2,5}) +4 \times 6 +3 \times 3  -8 =-25+24+9-8 = 0$, \quad le point de coordonnées $(\np{2,5}  ; 6 ; 3)$ appartient à (RST).

			Puisque le système \quad $\begin{cases}
				\np{2,5}=\np{-2,5} -10t \\
				6=8 + 4t \\
				3=\np{4,5} +3t
			\end{cases}$ \quad admet pour solution $t=-\dfrac12$, le point de coordonnées $(\np{2,5}  ; 6 ; 3)$ appartient à $d_1$.

			Le point de coordonnées $(\np{2,5}  ; 6 ; 3)$ appartient à $d_1$ et à (RST) (qui sont sécants) donc ce point est le point d'intersection de $d_1$ et (RST). C'est le point H, projeté orthogonal de A sur (RST).
		\end{enumerate}

		\item \begin{enumerate}
			\item $\vectt{SR} \cdot \vectt{ST}= 0 \times \np{2,5} -3 \times 4 +4 \times 3=0$

			Les vecteurs $\vectt{SR}$ et $\vectt{ST}$ sont orthogonaux donc les droites (SR) et (ST) sont orthogonale (et perpendiculaires, car sécantes en S) et le triangle RST est rectangle en S.

			\item Pour calculer le volume du tétraèdre, on choisit le triangle RST comme base et la hauteur associée est la droite (AH).

			Ainsi \quad $\mathcal{V}_{ARST} = \dfrac{1}{3} \times \dfrac{1}{2} SR \times ST \times AH$

			Or \quad  $SR=\sqrt{0^2+(-3)^2+(-4)^2}=5$,

			$ST=\sqrt{\np{2,5}^2+4^2+3^2}=\sqrt{\np{31,25}}$,

			et \quad $AH=\sqrt{5^2+(-2)^2+(-\np{1,5})^2}=\sqrt{\np{31,25}}$.

			Donc \quad $\mathcal{V}_{ARST} = \dfrac{1}{3} \times \dfrac{1}{2} \times 5 \times \sqrt{\np{31,25}} \times \sqrt{\np{31,25}}=\dfrac{625}{24}$.

			Comme l'unité graphique du repère est le centimètre, toutes les longueurs ont été exprimées en centimètre, et l'unité de volume est bien le cm$^3$, on a bien $\mathcal{V}_{ARST} = \dfrac{625}{24} \,\text{cm}^3$.
		\end{enumerate}
	\end{enumerate}

	\subsection*{Partie B}

	\begin{enumerate}
		\item $\begin{aligned}[t]
			\mathrm{AM}^2&=\left(x_M-x_A \right)^2+\left(y_M-y_A \right)^2+\left(z_M-z_A \right)^2
		=( -2+s)^2+(-4-4s)^2+(1-2s)^2\\
		&=4-4s+s^2+16+32s+16s^2+1-4s+4s^2
		=21s^2+24s+21
		\end{aligned}$

		On arrive bien à l'expression souhaitée.

		\item La fonction du second degré \quad $x \longmapsto 21s^2+24s+21$ \quad admet un minimum en \quad $\alpha = -\dfrac{b}{2a}=-\dfrac{4}{7}$.

		Ainsi la distance est minimale pour \quad $s=-\dfrac{4}{7}$ \quad soit pour \quad  M$\left(-\dfrac{71}{14} ; \dfrac{44}{7} ; \dfrac{61}{14} \right)$.
	\end{enumerate}

\end{document}
